\cotefile{157-2}
\scanpage{413}
\mspage{152}
seulement \unsure{en|sur|dans} $\mathrm{Hot}_{\unsure{0|o}}^{\wedge}$, ce qui n'est peut-être pas suffisant. Mais cela suffit pour \unsure{prouver} que la flèche $\unsure{\rho|\beta}$ \unsure{du} $(2)$ devient iso dans $\mathrm{Hot}_0^{\wedge}$, donc elle est iso.

$2^\circ)$ Il faut montrer que les HOT-fibrations
\diagram{\begin{tikzcd}[ampersand replacement=\&]
\deleted{\unsure{X}} \& \underline{\mathrm{Ch}}_{\infty}(X)
  \arrow[r,shift left=.5ex,"\sigma"] \arrow[r,shift right=.5ex,"\beta"'] \& X
\end{tikzcd}}
sont triviales, ou ce qui revient au \unsure{$=$|même}, \deleted{(puisqu'elles sont} qu'elles sont à fibres \deleted{triviales}\inserted{asphériques}. Il suffit de le voir pour $\sigma$, donc que
\[
\underline{\mathrm{Ch}}_{\infty}(X;a,\unsure{\underline{?}|?})
\]
est asphérique. Mais cette catégorie est la limite inductive des $\underline{\mathrm{Ch}}_n(X;a,?)$, i.e. des
\[
\underline{\mathrm{Hom}}^{\substack{\text{source}\\\text{fixée }a}}(\tau^n,\unsure{X}).
\]
Mais comme l'intervalle $I_{\unsure{\tau^n|\tau_n}}$ se rétracte par déformation sur son origine, il est \unsure{immédiat} que cette catégorie $\underline{\mathrm{Hom}}$ est même contractile.

$3^\circ)$ Considérons
\diagram{\begin{tikzcd}[ampersand replacement=\&,column sep=small]
 \& \overline{X'}=\deleted{\unsure{Z}}\unsure{E|\mathscr E|\mathbb E}(f)
 \arrow[dl,line width=1pt] \arrow[dr,"g"'] \arrow[ddrr,bend left=36,"\bar f"] \& \& \\
 X' \arrow[ur,dashed,"\alpha"] \arrow[dr,"f"'] \& |[font=\scriptsize]|{\underline{\text{cart}}} \&
 \underline{\mathrm{Ch}}_{\infty}(X) \arrow[dl,"\sigma"] \arrow[dr,dash,"\beta"'] \arrow[dr,shift left=1ex,dash,"\deleted{\unsure{X}}"] \& \\
 \& X \& \& X\,\deleted{\unsure{X}},
\end{tikzcd}}
il
% archive sheet number: 412
\scanpage{414}
\mspage{153}
\unsure{je voudrais|je voudrait} prouver$^{*}$ que $\bar f$ est une $\unsure{\mathbb D|\mathrm D}$-fibration, propre et lisse. \deleted{\unsure{$^{*}$|Ce}}
% slanted
\margin{\unsure{Le|Ce} résultat fondamental \unsure{de} \unsure{$^{*}$} \unsure{ce que} a été déjà vu, page 170-171}

$\overline{X'}$ est la catégorie des couples
\[
(x',c)\in\mathrm{Ob}\,X'\times\mathrm{Ch}_{\infty}(X)
\]
tels que
\[
\sigma(c)=f(x')
\]
et
\[
\bar f(x',c)=\beta(c)\qquad\text{but de }c.
\]

La fibre de $\deleted{\unsure{f}}\overline{X'},\bar f$ en $x_0\in X$ est donc \deleted{\unsure{formée}}\inserted{la cat.} des couples $(x',c)\in\mathrm{Ob}\,X'\times\mathrm{Ch}_{\infty}(X)$, tels que
\[
\sigma(c)=\unsure{f(x')},\quad\beta(c)=x_0.
\]

\deleted{\unsure{Cette}} On peut regarder encore $\overline{X'}$ comme limite inductive des sous-catégories $\overline{X'}_{n,m}$, \unsure{correspondant} aux couples $(x',c)$ avec $c\in\mathrm{Ch}_{n,m}(X)$. On trouve que pour $\unsure{m}$ \unsure{pair|paire} (direction \unsure{terminale|dominante} de $\tau_{n,m}$ rétrograde) les $\deleted{\unsure{\mathrm{Ch}}}\overline{X'}_{n,m}$ sont \inserted{Cat-}fibrées sur $X$, pour $m$ impair (direction \unsure{terminale|dominante} des $\tau_{n,m}$ directe) \deleted{elles} sont \inserted{Cat-}cofibrées. Donc on conclut encore comme tantôt, \deleted{par} $\overline{X'}$ étant \inserted{(au choix)} limite inductive \deleted{filtrante} de ces catégories lisses sur $X$, ou propres sur $X$,
% archive sheet number: 413
\scanpage{415}
\mspage{154}
Il faut prouver enfin que $\overline{X'}\mathrel{\overset{\bar f}{\unsure{\longrightarrow|\longleftrightarrow}}}\unsure{X'|X}$ est une $\unsure{\mathbb D|\mathrm D}$-fibration,$^{*}$ ce qui revient p. ex. à dire que pour une flèche
% slanted
\margin{\unsure{Cela|Celà} résulte déjà du fait qu'il est propre et lisse !}
\[
a\xrightarrow{\unsure{\alpha|u}}\deleted{\unsure{a}}b
\]
dans $X$, la flèche induite
\[
\overline{X'}_{/a}\xrightarrow{\bar{\unsure{v|u}}}\overline{X'}_{/b}
\]
est une Hot-équivalence. On sait déjà comme \unsure{tantôt} que l'on a un diagramme à flèches verticales des Hot-\unsure{équiv.|équi.}
\diagram{\begin{tikzcd}[ampersand replacement=\&]
\overline{X'}_{/a}\arrow[r] \& \overline{X'}_{/b}\\
\overline{X'}_a\arrow[u,hook,"\sim"']\arrow[r,dashed] \& \overline{X'}_b\arrow[u,hook,"\sim"'],
\end{tikzcd}}
et tout revient à montrer qu'une certaine flèche entre $\overline{X'}_a$ et $\overline{X'}_b$ dans Hot est \deleted{\unsure{iso}}. \deleted{\unsure{En}} J'ai envie de définir directement une telle flèche dans Hot, voire dans Cat. Or pour \inserted{des} $m$ \deleted{\unsure{est}} impairs, les $\overline{X'}_{n,m}$ sont Cat-cofibrées sur $X$. Pour $\overline{X'}_{-m,m}$ ($m$ impairs) on trouve donc
\[
(\overline{X'}_{-m,m})_a\xrightarrow{\unsure{u|v}_*}(\overline{X'}_{-m,m})_b
\]
% archive sheet number: 414
\scanpage{416}
\mspage{155}
\unsure{l'ennui|L'ennui} c'est qu'elles ne sont pas compatibles avec les morphismes de transition pour $m$ variable. Elles le sont \unsure{cependant} \uline{à} \uline{homotopie} \uline{près} \unsure{sûrement}, donc on trouvera au moins une flèche de Ind-objets dans Cat hot
\[
\overline{\overline{\deleted{\unsure{C}}X'}}_{\unsure{/a|a}}\xrightarrow{u_*}
\overline{\overline{X'}}_{\unsure{/b|b}}
\]
ce qui par passage à $\mathrm{Hot}_{\unsure{0|o}}^{\wedge}$ donne quand même une flèche
\[
\overline{X'}_{/a}\mathrel{\overset{u_*}{\unsure{\text{---}|\longrightarrow}}}\overline{X'}_{/b}
\]

En travaillant avec les $m$ \unsure{impairs}, on trouve de \unsure{$=$|même} des foncteurs en sens inverse
\[
(\overline{X'}_{-m,m})_a\xleftarrow{u^*}(\overline{X'}_{-m,m})_b,
\]
qui ne se recollent pas pour $m$ variable, mais qui se recollent à homotopie près, pour donner une flèche \deleted{\unsure{dans}} Ind Cat hot
\[
\overline{\overline{X'}}_{/a}\xleftarrow{\unsure{u^*|u_*}}\overline{\overline{X'}}_{/b}
\]
d'où aussi
\[
\overline{X'}_{a}\xleftarrow{u^*}\overline{X'}_{/b}\quad\text{dans }\mathrm{Hot}_0^{\wedge}.
\]
D'autre part, on doit \inserted{pouvoir} montrer \unsure{solidairement|solidement|successivement|soigneusement},
% archive sheet number: 415
\scanpage{417}
\mspage{156}
\unsure{cran par cran}, que $u_*u^*$ et $u^*u_*$ dans Ind Cat hot sont égales à l'identité, donc $u_*$ et $u^{\unsure{*|\wedge}}$ sont des iso. inverses \unsure{l'une|l'un} de l'autre. On doit vérifier d'autre part \unsure{strictement|solidement|solidairement} que
\diagram{\begin{tikzcd}[ampersand replacement=\&]
\overline{\overline{X'}}_{/a}\arrow[r,"\bar u"] \& \overline{\overline{X'}}_{/b}\\
\overline{\overline{X'}}_a\arrow[u,"\deleted{\unsure{\sim}}"]\arrow[r,"u_*"] \& \overline{\overline{X^{\unsure{''|'}}}}_b\arrow[u,"\deleted{\unsure{\sim}}"']
\end{tikzcd}}
dans Ind \deleted{\unsure{Hot}} \unsure{hCat|Hot} est commutatif.

Passant au \deleted{Ind Hot, $\mathrm{Hot}_{\unsure{0|o}}^\wedge$ et} diagramme correspondant dans $\mathrm{Hot}_0^{\wedge}$, en prenant les limites inductives, on trouve un carré commutatif dans $\mathrm{Hot}_0^{\wedge}$
\diagram{\begin{tikzcd}[ampersand replacement=\&]
\overline{X'}_{/a}\arrow[r,dash,"\bar u","{\unsure{\cdot|>}}"{pos=1,anchor=center}] \& \overline{X'}_{/b}\\
X'_a\arrow[u,"\sim"]\arrow[u,phantom,"\unsure{\xi|g}"']\arrow[r,"{\substack{u_*\\\sim}}"] \& X'_b\arrow[u,"\sim"],
\end{tikzcd}}
d'où résulte que $\bar u$ est iso dans $\mathrm{Hot}_0^{\wedge}$, donc aussi dans Hot.

Pour vérifier que la $\deleted{\unsure{\mathrm D}}\mathrm{HOT}$-fibration $\bar f$ est triviale \unsure{sss|ssi} $f$ est une \deleted{\unsure{équi}} HOT-équivalence, il faudrait avoir développé dans Cat le formalisme de la suite exacte \unsure{des|du|de} lacets -- on en est tout près\unsure{,} visiblement !
% archive sheet number: 416
\scanpage{418}
\mspage{157}
\uline{13}\quad Où on prouve enfin que $\underline{\mathrm{Ch}}(X,-,-)$ est aussi bon (voire meilleur) que $\underline{\mathrm{Ch}}_{\infty}(X,-,-)$.\hfill 13.11.90

$\cdot$ Je crois finalement que les catégories $\underline{\mathrm{Ch}}(X)$, $\underline{\mathrm{Ch}}(X;a,b)$ (sans \unsure{indices|indice} $\infty$ !) sont ``bonnes'', que les flèches \inserted{\unsure{can}} dans Hot
\[
\left\{\begin{aligned}
\underline{\mathrm{Ch}}_{\infty}(X)&\xrightarrow{\sim}\underline{\mathrm{Ch}}(X)\\
\underline{\mathrm{Ch}}_{\infty}(X;a,b)&\xrightarrow{\sim}\underline{\mathrm{Ch}}(X;a,b)
\end{aligned}\right.
\]
sont des iso. (La première est même iso dans $\mathrm{HOT}(X\times X)$, ce qui implique la deuxième relation, en passant aux fibres\ldots) Je crois que je suis à présent (après \unsure{XI|X1}) en mesure de le prouver. Cela résultera des deux résultats-clef suivants :

1) Le foncteur canonique
\[
X\longrightarrow\underline{\mathrm{Ch}}(X)
\]
(associant à \deleted{\unsure{tout}} $x\in\mathrm{Ob}\,X$ le chemin vide de source et but $x$) est une Hot-équivalence, et

2) Le foncteur source-but
\[
\underline{\mathrm{Ch}}(X)\longrightarrow X\times X
\]
est \uline{lisse} et \uline{propre} (donc une ``Hot-fibration parfaite'').

Cela doit donner ce qu'on cherche. Comme la \inserted{$X\times X$-}flèche $\underline{\mathrm{Ch}}_{\infty}(X)\mathrel{\unsure{\text{---}|\longrightarrow}}\underline{\mathrm{Ch}}(X)$ \inserted{(et les flèches déduites $\underline{\mathrm{Ch}}_\infty(X;a,b)\to\ldots$)} a été définie par l'intermédiaire \deleted{\unsure{des}} de la catégorie
\[
\mathcal C_{\mathbb N^{\unsure{o|0}}}=\mathbb N^o\times_{\underline{\unsure{\mathrm{Tz}|\mathrm{Ty}|\mathrm{Iz}|\mathrm{Iy}}}}\mathcal C
\]
\margin{(où $\mathcal C=\underline{\mathrm{Ch}}(X)$ considérée comme catégorie fibrée sur $\deleted{\unsure{\mathrm T}}\underline{\mathrm{Tz}}$ et $\mathbb N^o\to\underline{\mathrm{Tz}}$ est le syst. projectif des $\underbrace{(\Delta_+\Delta_-)^{\unsure{2n|2m}}}_{\unsure{t_n|\tau_n}}$ \ldots)}
% archive sheet number: 417
\scanpage{419}
\mspage{158}
à l'aide des deux flèches canoniques
\diagram{\begin{tikzcd}[ampersand replacement=\&]
\mathcal C_{\mathbb N^{\unsure{o|0}}}\arrow[r]\arrow[d] \& \underline{\mathrm{Ch}}_{\infty}(X)\\
\underline{\mathrm{Ch}}(X) \& {,}
\end{tikzcd}}
on sera obligé de prouver aussi les résultats qui correspondent à 1) 2) pour $\mathcal C_{\mathbb N^o}$, i.e.

% label at left of the diagram
\diagram{\begin{tikzcd}[ampersand replacement=\&,row sep=small]
{1')}\& X\arrow[r] \& \mathcal C_{\mathbb N^o}\\
\& \mathcal C_{t_0}=\unsure{C|\mathcal C}_0\arrow[u,phantom,"\simeq"{description,rotate=90}] \&
\end{tikzcd}}
(l'inclusion canonique de $C_0$, fibre de $\mathcal C_{\mathbb N^o}$ (fibré \unsure{en|sur} $\mathbb N^o$) en $0\in\mathbb N^o$) est une Hot-équivalence, et
\begingroup\makeatletter\tagsleft@true\makeatother
\[
\mathcal C_{\mathbb N^o}\longrightarrow X\times X\qquad\text{est propre et lisse.}\tag*{$2')$}
\]
\endgroup

Moyennant ces \deleted{\unsure{rés}} quatre résultats techniques, que je vais prouver tantôt (et $1'$ et $2'$ ne font aucune surprise pour moi), on gagne. Considérons en effet le diagramme
\diagram{\begin{tikzcd}[ampersand replacement=\&]
\&X\arrow[dl,"\alpha_1"']\arrow[d,"\alpha_0"]\arrow[dr,"\alpha_2"]\&\\
\underline{\mathrm{Ch}}(X)\arrow[dr,"p_1"] \& \mathcal C_{\mathbb N^o}\arrow[l,"\unsure{\pi|\Pi}_1"']\arrow[r,"\pi_2"]\arrow[d,"p_0"] \& \underline{\mathrm{Ch}}_\infty(X)\arrow[dl,"p_2"']\\
\&X\times X\&{.}
\end{tikzcd}}
Par 1) et $1')$ (concernant $\alpha_1$ et $\alpha_0$) et \unsure{le résultat|les résultats}
% archive sheet number: 418
\scanpage{420}
\mspage{159}
\unsure{semblable|similaire} déjà connu pour $\alpha_2$, ces trois foncteurs sont des Hot-équivalences. Donc il en est de $\unsure{=|\text{même}}$ \unsure{de|des} $\unsure{\pi|\Pi}_1,\pi_2$. Comme $p_1,p_0,p_2$ sont des Hot-fibrations (étant propres et lisses), il s'ensuit que $\pi_1$ et $\pi_2$ induisent des Hot équivalences sur les fibres
\diagram{\begin{tikzcd}[ampersand replacement=\&,column sep=large]
\underline{\mathrm{Ch}}(X;a,b) \& \mathcal C_{\mathbb N^{\unsure{o|0}}}(a,b)\arrow[l,"{\pi_1(a,b)}"']\arrow[r,"{\pi_2(a,b)}"] \& \underline{\mathrm{Ch}}_\infty(X;a,b).
\end{tikzcd}}
Je l'avais vérifié (sans mal) déjà pour $\pi_2$\inserted{ et les $\pi_2(a,b)$}, \unsure{ce} \deleted{\unsure{la chose de}}\inserted{qui avait} permis, en inversant $\pi_2$ dans Hot, de définir \uline{dans Hot}
\[
\left\{\begin{aligned}
\underline{\mathrm{Ch}}_\infty(X)&\mathrel{\unsure{\text{---}|\longrightarrow}}\underline{\mathrm{Ch}}(X)\\
\underline{\mathrm{Ch}}_\infty(X;a,b)&\longrightarrow\underline{\mathrm{Ch}}(X;a,b),
\end{aligned}\right.
\]
et dire que ce sont des iso signifie justement que $\pi_1$ et les $\pi_1(a,b)$ sont des iso.

Je ne crois pas que $\mathcal C_{\mathbb N^o}$ ait en lui-même d'autre intérêt, que de permettre de comparer les $X\times X$-catégories $\underline{\mathrm{Ch}}_\infty(X)$ et $\underline{\mathrm{Ch}}(X)$, à l'aide du diagramme précédent. Par contre, $\underline{\mathrm{Ch}}_\infty(X)$ et $\underline{\mathrm{Ch}}(X)$ m'apparaissent comme les deux ingrédients essentiels du formalisme des chemins dans \unsure{(Cat)|Cat}.
% archive sheet number: 419
\scanpage{421}
\mspage{160}
$\cdot$ Je vais prouver $1,2,1',2'$.

1) Considérons le diagramme
\diagram{\begin{tikzcd}[ampersand replacement=\&,column sep=small]
X=(X_{\unsure{\underline{\mathrm{Tz}}|\underline{\mathrm{Ty}}|\underline{\mathrm{Iz}}|\underline{\mathrm{Iy}}}})_e\arrow[r,hook,"i"]\arrow[dr,dashed] \&
X_{\underline{\mathrm{Tz}}}=X\times\underline{\mathrm{Tz}}\arrow[rr,"j"]\arrow[dr] \&\&\underline{\mathrm{Ch}}(X)\arrow[dl]\\
\& e\arrow[r,dashed] \&\underline{\mathrm{Tz}}\&{,}
\end{tikzcd}}
où $X_{\underline{\mathrm{Tz}}}$ et $\underline{\mathrm{Ch}}(X)$ sont considérées comme catégories Cat-fibrées sur $\underline{\mathrm{Tz}}$, et le foncteur $j$ est déduit des foncteurs canoniques
\[
X\longrightarrow\underline{\mathrm{Ch}}(X)_\tau\overset{\mathrm{def}}{=}\underline{\mathrm{Ch}}_\tau(X),
\]
associant à \deleted{chaque} $x\in\mathrm{Ob}\,X$ le chemin \deleted{vide} $\unsure{p|\varphi}_\tau^*(x)$, où $p_\tau:\tau\mathrel{\unsure{\text{---}|\to}}e$ est la projection de $\tau$ vers l'élément unité (et élément final) $e=1_{\underline{\mathrm{Tz}}}$ de $\underline{\mathrm{Tz}}$, et $\underline{\mathrm{Ch}}_e(X)$ est identifié à $X$. Il faut vérifier le

\uline{Lemme}\quad Pour tout $\tau\in\mathrm{Tz}$, le foncteur précédent est une Hot-équivalence (et même une \deleted{\unsure{équivale}} homotopisme).

Mais on sait que pour deux catégories $\unsure{\mathscr X|\mathcal X},\mathcal C$ lisses sur une \unsure{autre|même} $\unsure{T_0|T}$, tout $T$-morphisme de $\mathscr X$ dans $\mathcal C$ qui induit des Hot-équivalences sur les fibres, est une Hot-équivalence. Ainsi $j$ est
% archive sheet number: 420
\scanpage{422}
\mspage{161}
une Hot-équivalence. D'autre part, $\unsure{\underline{\mathrm{Tz}}|\underline{\mathrm{Ty}}|\underline{\mathrm{Iz}}|\underline{\mathrm{Iy}}}$ ayant un objet final est asphérique, donc $X\times\underline{\mathrm{Tz}}\xrightarrow{\mathrm{pr}_1}X$ est une Hot-équivalence, d'où résulte aussitôt que $i$ est une Hot-équivalence. Donc de même $ji$, or $ji=\alpha_1$, qed.

2) Prouvons p. ex. la lissité de $\underline{\mathrm{Ch}}(X)\to X\times X$, la propreté s'en déduira en \deleted{\unsure{utilisant}}\inserted{appliquant} ce résultat \deleted{pour}\inserted{à} $X^{\mathrm{opp}}$. Soit donc
\[
\unsure{\backslash}(a',b')\longrightarrow(\unsure{a},\unsure{b})\qquad\text{i.e.}\quad a'\mathrel{\overset{u}{\unsure{\text{---}|\to}}}a,\ b'\xrightarrow{v}b
\]
une flèche dans $X\times X$, et $c$ un objet de $\underline{\mathrm{Ch}}(X)$ au dessus de $(a,b)$ :
\diagram{\begin{tikzcd}[ampersand replacement=\&,column sep=huge]
a'\arrow[r,dashed,"c'","{\rightarrow}"']\arrow[d,"u"'] \& b'\arrow[d,"v"]\\
a\arrow[r,dash,"{\substack{\rightarrow\\\deleted{\unsure{\leftarrow}}\\c}}"'] \& b.
\end{tikzcd}}
On regarde la catégorie $\mathcal C(c,u,v)$ des couples $(c',w)$ qui \unsure{relèvent|relivent|relient} $(c,(u,v))$, \deleted{avec \unsure{$\mathcal C$ cat.}}\inserted{où} \deleted{\unsure{de|dans}} $c'$ est un chemin (de type non précisé) de $a'$ vers $b'$, et \deleted{\unsure{avec}}
\[
w:c'\longrightarrow c
\]
une flèche dans $\underline{\mathrm{Ch}}(X)$ qui sur les \inserted{objets-}sources induise $u$, et sur les objets buts induise $v$ (ainsi $c'$ est un chemin de $a'$ vers $b'$). Il faut
% archive sheet number: 421
\scanpage{423}
\mspage{162}
prouver que cette catégorie $\mathcal C(c,u,v)$ est asphérique. Pour cela, on regarde le chemin composé
\[
\widetilde c=\mathrm{ch}(u)\,c\,\mathrm{ch}(v)^{\unsure{o|0}}
\]
où $\mathrm{ch}(u)$ est le chemin \inserted{(direct)} de $a'$ \unsure{vers} $a$ défini par $u:a'\mathrel{\unsure{\text{---}|\to}}a$, et $\mathrm{ch}(\unsure{u|v})^o$ désigne donc le chemin (rétrograde) de $b$ à $b'$ défini par $v:b'\to b$. \deleted{\unsure{du}} $c$ est de type $\tau$, $\widetilde c$ est donc de type $\Delta_+\tau\Delta_-$.

On va comparer $\mathcal C(c;u,v)$ et la catégorie $\underline{\mathrm{Ch}}(X)_{/\widetilde c}$, à l'aide de deux foncteurs
\diagram{\begin{tikzcd}[ampersand replacement=\&,column sep=large]
\mathcal C(c,u,v)\arrow[r,shift left=.6ex,"\varphi"] \& \underline{\mathrm{Ch}}(X\unsure{,|;}a',b')_{/\widetilde c}\arrow[l,shift left=.6ex,"\psi"],
\end{tikzcd}}
dont on prouvera qu'ils définissent des flèches inverses l'une de l'autre dans Hot. Comme $\underline{\mathrm{Ch}}(X)_{/\widetilde c}$ est contractile (ayant un objet final), donc \deleted{\unsure{Hot}}asphérique, il en est de même de $\mathcal C(c,u,v)$, et on gagne.

\diagram{\begin{tikzcd}[ampersand replacement=\&,row sep=small,column sep=1em]
\&\&\makebox[0pt]{$\overbrace{\hspace{6cm}}^{\widetilde{c'}}$}\&\&\\
a'\arrow[r,"\mathrm{id}"]\arrow[d,"\mathrm{id}"'] \& a'\arrow[rr,dashed,dash,"c'"]\arrow[d,"u"'] \& {}\arrow[d,dashed,"w'"] \& b'\arrow[d,"v"'] \& b'\arrow[l,"\mathrm{id}"']\arrow[d,"\mathrm{id}"]\\
a'\arrow[r,"u"] \& a\arrow[rr,dash,"c"'] \& {} \& b \& b'\arrow[l,"v"']\\
\&\&\makebox[0pt]{$\underbrace{\hspace{6cm}}_{\widetilde c}$}\&\&
\arrow[from=2-1,to=3-2,phantom,"*" description]
\arrow[from=2-4,to=3-5,phantom,"*" description]
\end{tikzcd}}
On définit $\varphi$ par
\[
\boxed{\varphi(c',w')=(\widetilde{c'},\widetilde{w'})}
\]
où ($c'$ étant de type $\tau'$) $\widetilde{c'}$ est le chemin de type $\Delta^+\tau'\Delta^-$ déduit de \deleted{$\unsure{\varphi}$} $c'$ \inserted{\unsure{comme image inv.|comme image inverse|canoniquement}} par la dégénérescence $\Delta^+\tau'\Delta^-\to\tau'$, i.e. on prolonge $c'$ d'un ``cran \inserted{$\Delta_+^{\deleted{\unsure{+}}}$}'' $a'\xrightarrow{\mathrm{id}}a'$ à gauche, et d'un ``cran $\Delta_-^{\deleted{\unsure{-}}}$'' $b'\xleftarrow{\mathrm{id}}b'$ à droite, cf figure.
% archive sheet number: 422
\scanpage{424}
\mspage{164}
aussitôt que l'on a comm. \unsure{du|des}
\diagram{\begin{tikzcd}[ampersand replacement=\&]
\widetilde{c'}\arrow[rr,"\widetilde\alpha"]\arrow[dr,"\widetilde{w'}"'] \& \& \widetilde{c'_{\unsure{1|i}}}\arrow[dl,"\widetilde{w'_1}"]\\
\&\widetilde c\& \& {\text{i.e.}\quad\widetilde{w'}=\widetilde{w'_1}\widetilde\alpha}
\end{tikzcd}}
i.e. que $\widetilde\alpha$ est bien une flèche de $\underline{\mathrm{Ch}}(X,\deleted{\unsure{x}}a',b')_{/\widetilde c}$. De $=$, que $\alpha\mapsto\widetilde\alpha$ définit une dépendance fonctorielle, i.e. $\widetilde{\beta\alpha}=\widetilde\beta\,\unsure{\widehat\alpha|\widetilde\alpha}$, est évident sur la définition de $\widetilde\alpha$. \deleted{\unsure{Cela}} Donc, on a bien défini
\[
\mathcal C(c,u,v)\xrightarrow{\varphi}\underline{\mathrm{Ch}}(X;a',b')_{/\widetilde c},
\]
définissons maintenant
\[
C(c,u,v)\xleftarrow{\psi}\underline{\mathrm{Ch}}(X;a',b')_{/\widetilde c}.
\]
Soit $c'',w''$ un objet de $\underline{\mathrm{Ch}}(X;a',b')_{/\widetilde c}$,
\diagram{\begin{tikzcd}[ampersand replacement=\&,row sep=small,column sep=1em]
a'_0\arrow[dr,equals]\&\&\&\makebox[0pt]{$\overbrace{\hspace{6.2cm}}^{c''}$}\&\&\\
\& a'\arrow[r]\arrow[d,dash,"\mathrm{id}"'] \& a'_1\arrow[rr,dash,"\rightarrow"]\arrow[d,"w''_1"] \& {}\arrow[d,dashed,"w''"] \& a'_n\arrow[d,"w''_n"'] \& b'\arrow[l]\arrow[d,equals]\\
\& a'\arrow[r,"u"] \& a\arrow[rr,dash] \& {} \& b \& b'\arrow[l,"v"'],\\
\&\&\&\makebox[0pt]{$\underbrace{\hspace{3.3cm}}_{c}$}\&\&
\end{tikzcd}}
donc
\[
w'':c''\mathrel{\unsure{\text{---}|\longrightarrow}}\widetilde c\qquad\text{flèche de }\underline{\mathrm{Ch}}(X;a',b');
\]
notons qu'on a une flèche de rétraction évidente
\[
\widetilde c\xrightarrow{\rho_c}c\deleted{\unsure{=\varphi}\,\unsure{\delta_{\Delta_+}}\unsure{\tau}\,\unsure{\delta_{\Delta_-}}}
\]
\inserted{associée à $\delta_{\tau}:\Delta_+\tau\Delta_-\to\tau$ (rétraction)} \unsure{induisant|traduisant} $u$ et $v$ sur la source et le but :
\diagram{\begin{tikzcd}[ampersand replacement=\&,column sep=.55em,row sep=small]
\deleted{\unsure{a}}\&a'\arrow[r,"{\overset{\frown}{u}}"]\arrow[d,"u"] \& a\arrow[rrr,dash,"c"]\arrow[d,dash,"\mathrm{id}"]\&{}\arrow[d,"\mathrm{id}"']\&{}\arrow[d,"\mathrm{id}"]\& b\arrow[d,"\mathrm{id}"]\&b'\arrow[l,"v"']\arrow[d,"v"]\&\big\}\ \widetilde c\in\mathrm{Ch}_{\Delta_+\tau\Delta_-}(a',b',\\
\&a\arrow[r,"\mathrm{id}"]\&a\arrow[rrr,dash,"c"]\&\&\&b\&b\arrow[l,"\mathrm{id}"']\&\big\}\ \delta_\tau^*(c)\in\mathrm{Ch}_{\Delta_+\tau\Delta_-}(a,b)\\
\&\&a\arrow[rrr,dash,"c","{\leftarrow}"']\&\&\&b\&\&\big\}\ c\in\mathrm{Ch}_\tau(a,b).
\end{tikzcd}}
% archive sheet number: 423
\scanpage{425}
\mspage{163}
On peut dire que
\[
\widetilde{c'}=a'_{\Delta_+}\unsure{\circ|*}\unsure{\tau'|c'}\circ b'_{\Delta_-},
\]
où de façon générale, si $x\in X$ et $\tau\in\unsure{\mathrm{Tz}|\mathrm{Ty}|\mathrm{Iz}|\mathrm{Iy}}$, $x_\tau$ désigne le ``chemin constant'' de valeur $x$, type $\tau$.

On a d'autre part des morphismes évidents
\[
a'_{\Delta_+}\longrightarrow\mathrm{ch}(u),\qquad b'_{\Delta_-}\longrightarrow\mathrm{ch}(v)^{\mathrm{opp}},
\]
\unsure{comme ci-contre|comme montré} par les flèches verticales $\mathrm{id}_{a'},u$ et $v,\mathrm{id}_{b'}$ dans les deux carrés \inserted{marqués} $*$ du diagramme, flèches qui sont l'une dans $\underline{\mathrm{Ch}}_{\Delta_+}(X)$, l'autre dans $\underline{\mathrm{Ch}}_{\Delta_-}(X)$. Ceci dit, on pose
\[
\widetilde{w'}:\widetilde{c'}\mathrel{\unsure{\text{---}|\longrightarrow}}\widetilde c,\qquad
\widetilde w=\deleted{\unsure{\mathrm{ch}}}\lambda(u)*w'*\deleted{\unsure{\mathrm{ch}}}\mu(v)^{\deleted{\unsure{\mathrm{opp}}}}
\]
où
\[
\left\{\begin{aligned}
\lambda(u)&:a'_{\Delta_+}\longrightarrow\mathrm{ch}(u)\\
\mu(v)&:b'_{\Delta_-}\longrightarrow\mathrm{ch}(v)^{\mathrm{opp}}
\end{aligned}\right.
\]
(cf figure). Le couple
\[
(\widetilde{c'},\widetilde{w'})\in\mathrm{Ob}\,\underline{\mathrm{Ch}}(X)_{/\widetilde c}
\]
ainsi défini dépend fonctoriellement de $(\deleted{\unsure{\widetilde c}}c',w')$ variable dans $\mathcal C(c,u,v)$ : si $(c'_1,w'_1)$ est un deuxième objet de cette catégorie, une flèche du premier dans le second est une flèche \unsure{$\alpha$}
\[
c'\xrightarrow{\alpha}c'_1
\]
dans $\underline{\mathrm{Ch}}(X;a',b')$, telle qu'on ait comm. du
\diagram{\begin{tikzcd}[ampersand replacement=\&]
c'\arrow[rr,dash,"\alpha","{\unsure{\cdot|>}}"{pos=1,anchor=center}]\arrow[dr,dash,"w'"']\&\&c'_1\arrow[dl,dash,"w'_1"]\& {w'=w'_1\alpha}\\
\&c\&{.}\&
\end{tikzcd}}
Une telle flèche $\alpha$ définit
\[
\widetilde\alpha=a'_{\Delta_+}*\alpha*b'_{\Delta_-}:\widetilde{c'}\longrightarrow\widetilde{c'_1},
\]
\inserted{dans $\underline{\mathrm{Ch}}(X,a',b')$} et on vérifie
% archive sheet number: 424
\scanpage{426}
\mspage{165}
\deleted{\illegible{1}} Donc on définit un objet \deleted{de $C(c,u,v)$ \unsure{$\mid$}}
\[
(c',w')=\psi(c'',w'')\in\mathrm{Ob}\,C(c,u,v)
\]
par
\[
c^{\unsure{\prime|\prime\prime}}=c'',\qquad w'=\rho_c w'',
\]
et c'est fonctoriel en $(c'',w'')$.

Le composé $\varphi\psi$ étant un endomorphisme de la catégorie contractile $\underline{\mathrm{Ch}}(a',b')_{/\widetilde c}$, est homotope à l'identité. Il faut vérifier qu'il en est de $=$ de $\psi\varphi$, endomorphisme de $C(c,u,v)$. Ce foncteur associe à $(c',w')$, $(\widetilde{c'},\overline{w'})$, où $\widetilde{c'}$ est défini comme plus haut $\widetilde{c'}=a'_{\Delta_+}\circ c'\circ b'_{\Delta_-}$, et où $\overline{w'}$ est le composé
\diagram{\begin{tikzcd}[ampersand replacement=\&,column sep=small,row sep=small]
\widetilde{c'}\arrow[rr,"\rho_{c'}",""{name=rhoanchor,below}]\arrow[d,equals] \&\& c'\arrow[r,"w'"]\&c\\
\delta_{\tau'}^{\unsure{*|'*}}(c') \& \& {\substack{\text{\unsure{flèche}}\\\text{canonique}}} \arrow[from=rhoanchor,to=2-3,dash,bend left=12] \&
\end{tikzcd}}

On a donc
\[
(\psi\varphi)(c',w')=(\widetilde{c'},\overline{w'})\xrightarrow{\rho_{c'}}(c',w')
\]
c'est bien un morphisme dans $C(c,u,v)$ i.e. on a commutativité \unsure{du|des}
\diagram{\begin{tikzcd}[ampersand replacement=\&]
\widetilde{c'}\arrow[rr,"\rho_{c'}"]\arrow[dr,"\overline{w'}"'] \&\& c'\arrow[dl,"w'"]\&{\text{(flèche de }\mathrm{Ch}(a',b')}\\
\&c\&\&\unsure{\text{--}}
\end{tikzcd}}
% archive sheet number: 425
\scanpage{427}
\mspage{166}
Ainsi, le foncteur $\psi\varphi$ est bien homotope au foncteur identique. On a gagné !

Passons à $1')$, $2')$ pour $\mathcal C_{\mathbb N^{\unsure{o|0}}}$. L'assertion $1')$ se démontre exactement de la $=$ manière que 1) pour $\underline{\mathrm{Ch}}(X)$, en regardant $\mathcal C_{\mathbb N^o}$ comme fibrée sur $\deleted{\underline{\unsure{\mathrm{Tz}|\mathrm{Ty}|\mathrm{Iz}|\mathrm{Iy}}}}\mathbb N^o$ (qui a un objet final $0$), et en \deleted{\unsure{du}}utilisant le diagramme
\diagram{\begin{tikzcd}[ampersand replacement=\&,column sep=small]
X\arrow[r,"i'"]\arrow[d,equals] \& X_{\deleted{\underline{\mathrm{Tz}}}\mathbb N^o}=X\times\deleted{\underline{\mathrm{Tz}}}\mathbb N^o\arrow[rr,"j'"]\arrow[dr]\&\&\mathcal C_{\mathbb N^o}\arrow[dl]\\
\deleted{\unsure{=}}(X_{\deleted{\underline{\mathrm{Tz}}}\mathbb N^o})_{\unsure{e|0}}\&\&\deleted{\underline{\mathrm{Tz}}}\mathbb N^o\&{,}
\end{tikzcd}}
\deleted{\unsure{où}} où $j'$ est (cartésien et) une Hot-équivalence sur chaque fibre, d'après le lemme, donc une Hot équivalence \deleted{\unsure{part}} globalement, et $i'$ est une Hot-équivalence car $\mathbb N^o$ est asphérique (ayant un objet final). Donc $j'i'$ est asphérique qed.

Il reste à prouver que
\[
C_{\mathbb N^o}\mathrel{\unsure{\text{---}|\longrightarrow}}X\times X
\]
est propre et lisse. \deleted{Par \unsure{dualité}, \unsure{passage|passant} de $X$ à $X^{\unsure{o|'}}$, il suffit de prouver que c'est lisse\unsure{.|,}} Ici encore,
% archive sheet number: 426
